\chapter{Martyngały w czasie ciągłym: informacja, regularność i lokalizacja}
\label{ch:mart-cont}

\section{Dlaczego przejście od czasu dyskretnego nie jest automatyczne?}
W tomie I (rozdział~\ref{I-ch:nadmartyngaly})
badaliśmy martyngały w czasie dyskretnym.
Martyngał był ciągiem $(M_n)_{n\ge0}$, którego
warunkowa prognoza następnej wartości, przy znajomości przeszłości,
wynosiła dokładnie wartość obecną. W czasie ciągłym rozważamy
nieprzeliczalną rodzinę chwil. Pojawiają się dwa nowe pytania:
jak wybierać jedną wspólną wersję trajektorii oraz co znaczy
zatrzymać proces w chwili wybranej na podstawie obserwacji?
Później zobaczymy też, że proces ze skokami może być martyngałem,
mimo że jego trajektorie mają skończoną wariację całkowitą.

\begin{definicja}[Martyngał, podmartyngał i nadmartyngał]
Niech $(\Omega,\calF,\Prob)$ będzie przestrzenią probabilistyczną,
$(\calF_t)_{t\ge0}$ filtracją i $M=(M_t)_{t\ge0}$ procesem
rzeczywistym, adaptowanym, takim że $\E|M_t|<\infty$
dla każdego $t\ge0$. Proces $M$ jest \emph{martyngałem}, jeśli
$\E[M_t\mid\calF_s]=M_s$ prawie na pewno dla wszystkich
$0\le s\le t$. Jest \emph{podmartyngałem}, jeśli
$\E[M_t\mid\calF_s]\ge M_s$, a \emph{nadmartyngałem},
gdy zachodzi nierówność przeciwna. Filtracja spełnia
\emph{warunki zwykłe}, jeśli jest zupełna (każdy podzbiór zbioru z $\calF$ o prawdopodobieństwie
zero należy do $\calF_0$) i prawostronnie ciągła:
$\calF_t=\bigcap_{u>t}\calF_u$ dla każdego $t\ge0$.
\end{definicja}

Ważne, że równość dla każdej pary $s,t$ jest rozumiana osobno
prawie na pewno. Nie wolno bez dodatkowego argumentu zażądać
jednego zbioru miary jeden, na którym równości warunkowe
miałyby zachodzić jednocześnie dla nieprzeliczalnie wielu par.

\begin{przyklad}[Trzy modele uczciwej ewolucji]
Niech $W$ będzie standardowym procesem Wienera w jego
uzupełnionej filtracji naturalnej. Wówczas
$W_t$ jest martyngałem z powodu niezależnych wycentrowanych
przyrostów (rozdział~\ref{ch:wiener}). Uzupełnienie filtracji
o zbiory zerowe nie zmienia równości warunkowych. Proces $W_t^2-t$ również nim jest, ponieważ
\[
\E[W_t^2-t\mid\calF_s]
=W_s^2+\E[(W_t-W_s)^2]-t=W_s^2-s.
\]
Jeśli $N$ jest procesem Poissona o intensywności $\lambda>0$
w swojej filtracji naturalnej (rozdział~\ref{ch:poisson-process}),
to $N_t-\lambda t$ jest martyngałem, gdyż przyrost liczby
zdarzeń po chwili $s$ jest niezależny od przeszłości i ma
średnią $\lambda(t-s)$. Pierwsze dwa procesy mają trajektorie
ciągłe, ostatni ma skoki wysokości jeden.
\end{przyklad}

\begin{lemat}[Warunkowe wartości oczekiwane są jednostajnie całkowalne]
\label{lem:mart-cont-ui-cond}
Niech $X\in L^1(\Omega,\calF,\Prob)$. Rodzina
$\{\E[X\mid\calG]:\calG\subseteq\calF\text{ jest pod-}\sigma\text{-ciałem}\}$
jest jednostajnie całkowalna.
\end{lemat}
\begin{proof}
Ustalmy $K>0$ i rozłóżmy $X=U+V$, gdzie
$U=X\chi_{\{|X|\le K/2\}}$ oraz
$V=X\chi_{\{|X|>K/2\}}$. Dla dowolnego $\calG$
połóżmy $Y=\E[X\mid\calG]$. Wówczas
$|\E[U\mid\calG]|\le K/2$, skąd na zdarzeniu
$\{|Y|>K\}$ zachodzi $|\E[V\mid\calG]|>K/2$
i $|Y|\le K/2+|\E[V\mid\calG]|\le2|\E[V\mid\calG]|$.
Dlatego
\[
\E[|Y|\chi_{\{|Y|>K\}}]
\le2\E|\E[V\mid\calG]|
\le2\E|V|
=2\E[|X|\chi_{\{|X|>K/2\}}].
\]
Prawa strona nie zależy od $\calG$ i dąży do zera
przy $K\to\infty$ dzięki całkowalności $X$.
\end{proof}

\begin{wniosek}[Jednostajna całkowalność na skończonym horyzoncie]
Niech $M$ będzie martyngałem oraz $T<\infty$.
Wtedy rodzina $\{M_t:0\le t\le T\}$ jest jednostajnie
całkowalna, gdyż $M_t=\E[M_T\mid\calF_t]$.
\end{wniosek}

\section{Prawostronne modyfikacje: skąd bierze się regularność?}
\begin{definicja}[Trajektoria c\`adl\`ag]
Funkcja $f:[0,\infty)\to\R$ jest c\`adl\`ag, jeśli
jest prawostronnie ciągła i ma skończone granice lewostronne
we wszystkich chwilach $t>0$. Proces jest c\`adl\`ag,
gdy prawie wszystkie jego trajektorie mają tę własność.
\end{definicja}

\begin{lemat}[Zbieżność odwrotnych warunkowych wartości oczekiwanych]
Niech $(\calG_n)_{n\ge1}$ będzie malejącym ciągiem
pod-$\sigma$-ciał przestrzeni $(\Omega,\calF,\Prob)$,
$\calG_\infty=\bigcap_n\calG_n$, a $X\in L^1$.
Wtedy
$\E[X\mid\calG_n]\to\E[X\mid\calG_\infty]$ w $L^1$.
\end{lemat}
\begin{proof}
Najpierw niech $X\in L^2$ i $Y_n=\E[X\mid\calG_n]$.
Dla $m\ge n$ mamy $Y_m=\E[Y_n\mid\calG_m]$.
Ortogonalność rzutu warunkowego daje
$\|Y_n-Y_m\|_2^2=\|Y_n\|_2^2-\|Y_m\|_2^2$.
Normy $\|Y_n\|_2$ maleją i są nieujemne, zatem
$(Y_n)$ jest ciągiem Cauchy'ego w $L^2$. Niech $Y$
będzie jego granicą. Wybierzmy podciąg $(Y_{n_j})$ zbieżny prawie na pewno
do $Y$, co jest możliwe dzięki zbieżności w $L^2$.
Aby uzyskać jedną wersję mierzalną względem
przecięcia także bez założenia zupełności $\calG_k$,
definiujemy $\widehat Y$ jako granicę tego podciągu
tam, gdzie istnieje skończona granica, a poza tym
jako zero. Dla każdego $k$ definicja zależy wyłącznie
od ogona podciągu o indeksach $n_j\ge k$, dlatego
$\widehat Y$ jest $\calG_k$-mierzalna dla wszystkich $k$
i w konsekwencji $\calG_\infty$-mierzalna.
Ponadto $\widehat Y=Y$ p.n.; dalej używamy tej wersji.
Dla $A\in\calG_\infty$ i wszystkich $n$ mamy
$\E[Y_n\chi_A]=\E[X\chi_A]$; przejście do granicy
w $L^1$ identyfikuje $Y=\E[X\mid\calG_\infty]$.
Dla ogólnego $X\in L^1$ przybliżamy $X$ obcięciami
$X^{(K)}\in L^2$. Kontrakcja warunkowania w $L^1$
pozwala przenieść zbieżność z $X^{(K)}$ na $X$ po
przejściu $K\to\infty$.
\end{proof}

\begin{lemat}[Dlaczego z ograniczonej liczby przekroczeń wynikają granice?]
Niech $D$ będzie gęstym podzbiorem przedziału $[0,T]$,
a $f:D\to\R$ funkcją ograniczoną. Załóżmy, że dla każdej
pary liczb wymiernych $a<b$ funkcja $f$, oglądana w
rosnącym porządku chwil z $D$, ma skończenie wiele pełnych
przekroczeń przedziału $[a,b]$ w górę. Wtedy w każdej chwili
$t<T$ istnieje skończona prawa granica
$\lim_{r\downarrow t,\,r\in D}f(r)$, a w każdej chwili
$t>0$ istnieje skończona lewa granica wzdłuż $D$.
Funkcja zdefiniowana przez prawe granice (z wartością $f(T)$
w końcu przedziału, gdy $T\in D$) jest c\`adl\`ag.
\end{lemat}
\begin{proof}
Brak prawej granicy w chwili $t$ oznaczałby
$\liminf_{r\downarrow t,\,r\in D}f(r)
<\limsup_{r\downarrow t,\,r\in D}f(r)$.
Można wtedy wybrać wymierne $a<b$ pomiędzy tymi
liczbami oraz dowolnie długie skończone ciągi chwil
z $D\cap(t,t+\delta)$, w których wartości naprzemiennie
spadają poniżej $a$ i rosną powyżej $b$.
Po ustawieniu chwil rosnąco otrzymujemy dowolnie wiele
przekroczeń w górę, sprzecznie z założeniem.
To samo rozumowanie, z chwilami zbliżającymi się do $t$
od lewej, daje istnienie granicy lewostronnej.
Ograniczoność $f$ zapewnia skończoność obu granic.

Oznaczmy prawą granicę przez $g(t)$. Gdyby dla
$t_n\downarrow t$ zachodziło $g(t_n)\not\to g(t)$,
można by znaleźć $r_n\in D$ położone tuż za $t_n$,
tak by $r_n\downarrow t$ oraz
$|f(r_n)-g(t_n)|\to0$; byłaby to sprzeczność z
definicją $g(t)$. Analogicznie, dla $t_n\uparrow t$
wybieramy $r_n\in D$ między $t_n$ i $t$.
Istnienie lewej granicy $f$ w $t$ pociąga wówczas
istnienie tej samej lewej granicy dla $g$.
Zatem $g$ jest prawostronnie ciągła i ma granice
lewostronne.
\end{proof}

\begin{twierdzenie}[Regularizacja martyngału typu zamkniętego]
\label{thm:regularizacja}
Niech filtracja $(\calF_t)_{t\ge0}$ spełnia warunki zwykłe,
a $X\in L^1(\Omega,\calF,\Prob)$. Proces
$M_t=\E[X\mid\calF_t]$ ma modyfikację c\`adl\`ag.
Dokładniej, istnieje proces $\widetilde M$ o trajektoriach
c\`adl\`ag, taki że dla każdego $t\ge0$
$\widetilde M_t=M_t$ prawie na pewno.
\end{twierdzenie}
\begin{proof}
Ustalmy skończony horyzont $T>0$ i wybierzmy
$D_T=(\Q\cap[0,T])\cup\{T\}$.
Wybieramy rosnący ciąg skończonych podzbiorów $D_T$
zawierających $0,T$ i wyczerpujących $D_T$.
Dla dowolnych liczb
wymiernych $a<b$ i skończonego uporządkowanego podzbioru
$D_T$ zastosujmy nierówność przekroczeń do dyskretnego
martyngału otrzymanego przez próbkowanie $(M_t)$ w tych
chwilach. Oczekiwana liczba przekroczeń przedziału
$[a,b]$ jest ograniczona przez stałą zależną od
$a,b$ i $\E|X|$, a nie od liczby punktów próbkowania:
taki wniosek wynika z nierówności przekroczeń
Dooba z tomu I, rozdział~\ref{I-ch:nadmartyngaly}:
dla martyngału na skończonej siatce
$(b-a)\E U[a,b]\le\E[(M_T-a)^-]$
i prawa strona jest nie większa od $\E|X|+|a|$.
Tutaj $U[a,b]$ oznacza liczbę pełnych przekroczeń
w górę, a $y^-=\max(-y,0)$.
Przy dodawaniu punktów siatki liczba przekroczeń
nie maleje (definiujemy ją jako maksymalną
liczbę rozłącznych kolejnych par punktów).
Zbieżność monotoniczna zapewnia skończoną liczbę
przekroczeń wzdłuż $D_T$ prawie na pewno.
Ponadto $|M|$ jest podmartyngałem przez nierówność
Jensena. Słaba nierówność maksymalna Dooba
na każdej skończonej siatce daje dla $K>0$
$\Prob(\max |M_r|>K)\le \E|M_T|/K\le\E|X|/K$.
Przechodząc do całego $D_T$, a następnie z $K$
do nieskończoności, uzyskujemy
$\sup_{r\in D_T}|M_r|<\infty$ prawie na pewno.
Na wspólnym zbiorze pełnej miary, wyznaczonym
przez te warunki dla przeliczalnie wielu par
wymiernych $a<b$, stosujemy poprzedni lemat.
Dostajemy skończone prawe granice dla $0\le s<T$
i lewe dla $0<s\le T$.

Na tym zbiorze definiujemy
$\widetilde M_s=\lim_{r\downarrow s,r\in D_T}M_r$
dla $s<T$ oraz $\widetilde M_T=M_T$.
Na jego dopełnieniu przyjmujemy zerową trajektorię. Poprzedni lemat
o funkcjach rzeczywistych z granicami jednostronnymi
zapewnia, że $s\mapsto\widetilde M_s$ jest prawostronnie
ciągła i ma granice z lewej: gdyby prawostronna ciągłość
zawiodła, powstałyby dwie sekwencje punktów zbieżne do
tej samej chwili o różnych granicach, sprzeczne z istnieniem
wspólnej granicy prawostronnej próbek z $D_T$.

Aby sprawdzić, że to modyfikacja właściwego procesu, dla
każdego ustalonego $s<T$ wybieramy $r_n\in D_T$ malejące
do $s$. Z lematu o malejących $\sigma$-ciałach i prawostronnej
ciągłości filtracji
$M_{r_n}=\E[X\mid\calF_{r_n}]\to
\E[X\mid\calF_s]=M_s$ w $L^1$.
Z drugiej strony lewa strona zbiega prawie na pewno do
$\widetilde M_s$. Stąd $\widetilde M_s=M_s$ prawie na pewno.
Dla ustalonego $s$ granica wzdłuż malejących
$r_n$ jest $\bigcap_n\calF_{r_n}=\calF_s$-mierzalna:
można ją zdefiniować jako granicę ogona tam,
gdzie istnieje skończona granica, i zero w przeciwnym
razie. Zmiana na wspólnym zbiorze zerowym jest
$\calF_0$-mierzalna dzięki zupełności filtracji.
Zatem $\widetilde M$ jest adaptowany.

Łącząc konstrukcje na horyzontach $T=1,2,\ldots$,
bierzemy wspólny zbiór pełnej miary, na którym
istnieją wszystkie potrzebne granice oraz
w każdym całkowitym punkcie końcowym wartość
$M_T$ jest zgodna z prawą granicą uzyskaną
na dłuższym horyzoncie. Ten ostatni warunek
wynika z udowodnionej własności modyfikacji
i przeliczalności całkowitych $T$.
W punktach wewnętrznych zgodność jest bezpośrednia,
bo używamy tych samych wymiernych chwil.
Poza wspólnym zbiorem pełnej miary definiujemy
zerową trajektorię. Otrzymujemy adaptowany proces
o wszędzie c\`adl\`ag trajektoriach, który jest
modyfikacją $M$, a przez równości w ustalonych
chwilach również martyngałem.
Każde dwie prawostronnie ciągłe modyfikacje
są nierozróżnialne, gdyż zgadzają się p.n.
w każdej z przeliczalnie wielu chwil wymiernych,
a prawostronna ciągłość rozszerza zgodność
na wszystkie chwile jednocześnie.
\end{proof}

\begin{twierdzenie}[Ogólne twierdzenie regularizacyjne Dooba]
\label{thm:doob-general-regularization}
Niech $(\calF_t)$ spełnia warunki zwykłe, a $X$ będzie
rzeczywistym podmartyngałem. Istnieje adaptowana
modyfikacja c\`adl\`ag wtedy i tylko wtedy, gdy
$t\mapsto\E X_t$ jest prawostronnie ciągła.
Analogiczna teza zachodzi dla nadmartyngałów.
W szczególności każdy martyngał ma taką modyfikację.
\end{twierdzenie}
\begin{proof}
Najpierw wykażemy istnienie granic trajektorii
na chwilach wymiernych, niezależnie od założenia
o ciągłości średniej. Dla $T>0$ połóżmy
$D_T=(\Q\cap[0,T])\cup\{T\}$.
Na każdej skończonej siatce z końcem $T$
nierówność przekroczeń Dooba (tom I,
rozdział~\ref{I-ch:nadmartyngaly}) daje
$(b-a)\E U[a,b]\le \E[(X_T-a)^+]+
\E[(X_0-a)^-]$ (wystarcza to oszacowanie).
Liczba przekroczeń wzdłuż całego $D_T$
jest więc skończona p.n. dla każdej wymiernej
pary $a<b$.

Aby uzyskać także ograniczoność próbek, połóżmy
$C_u=\E[X_T\mid\calF_u]$ i $Y_u=C_u-X_u$
dla $0\le u\le T$. Z własności podmartyngału
$Y_u\ge0$ p.n., a $Y$ jest nadmartyngałem.
Na skończonej siatce maksymalna nierówność
dla nieujemnego nadmartyngału daje
$\Prob(\max Y_u>K)\le\E Y_0/K$:
zatrzymujemy na pierwszym przekroczeniu poziomu
$K$ (lub końcu siatki), stosujemy dyskretne
próbkowanie opcjonalne i nieujemność.
Nierówność maksymalna dla $|C|$ daje analogiczne
oszacowanie przez $\E|X_T|/K$.
Wyczerpując $D_T$ skończonymi siatkami, otrzymujemy
$\sup_{u\in D_T}|X_u|<\infty$ p.n.
Lemat o przekroczeniach pozwala zatem zdefiniować
prawostronne granice $\widetilde X_u$ wzdłuż $D_T$
dla $u<T$, z trajektoriami c\`adl\`ag po dołączeniu
wartości $X_T$ na końcu. Jak w dowodzie
twierdzenia~\ref{thm:regularizacja}, warunki
zwykłe zapewniają adaptację tej konstrukcji.

Potrzebujemy teraz zbieżności w $L^1$ przy
$u_j\downarrow u$. Oznaczmy $V_j=X_{u_j}$,
$\mathcal G_j=\calF_{u_j}$ i
$A_j=\E[V_j\mid\mathcal G_{j+1}]-V_{j+1}\ge0$
p.n. Z własności wieży dostajemy indukcyjnie
\[
V_n=\E\left[V_1-\sum_{j=1}^{n-1}A_j
\,\middle|\,\mathcal G_n\right].
\]
Istotnie $A_j$ jest $\mathcal G_{j+1}$-mierzalna,
a $\mathcal G_{j+1}\subseteq\mathcal G_j$;
kolejne warunkowanie powyższej równości i
odjęcie $A_n$ daje następny krok.
Ponieważ $\E V_n\ge\E X_u$, mamy
\[
\sum_{j\ge1}\E A_j
=\E V_1-\lim_n\E V_n
\le \E V_1-\E X_u<\infty.
\]
Zatem $A=\sum_{j\ge1}A_j$ jest całkowalna
i $V_n$ różni się od
$\E[V_1-A\mid\mathcal G_n]$ w normie $L^1$
o co najwyżej $\E\sum_{j\ge n}A_j\to0$.
Lemat o odwrotnych warunkowych wartościach
oczekiwanych daje zbieżność $V_n$ w $L^1$.
Jeżeli $u_j\in D_T$, jej granica jest
$\widetilde X_u$, ponieważ istnieje granica p.n.

W szczególności $\E\widetilde X_u=
\lim_{j\to\infty}\E X_{u_j}$ oraz
$\E[V_j\mid\calF_u]\ge X_u$.
Przechodząc do granicy w $L^1$ i korzystając
z $\calF_u$-mierzalności $\widetilde X_u$,
dostajemy $\widetilde X_u\ge X_u$ p.n.
Gdy średnia jest prawostronnie ciągła,
różnica ta ma średnią zero, więc
$\widetilde X_u=X_u$ p.n.
Łączenie horyzontów i wybór zerowej trajektorii
na wspólnym zbiorze zerowym odbywa się
jak w twierdzeniu~\ref{thm:regularizacja}.

Odwrotnie, jeśli istnieje modyfikacja
prawostronnie ciągła, dla każdego ciągu
$u_j\downarrow u$ mamy $X_{u_j}\to X_u$ p.n.
(zgodność wersji wystarczy w przeliczalnym
zbiorze tych chwil). Udowodniona powyżej
zbieżność w $L^1$ dla takiego ciągu pokazuje
$\E X_{u_j}\to\E X_u$. Ponieważ średnia
podmartyngału jest niemalejąca, to oznacza
jej prawostronną ciągłość.
Dla nadmartyngałów stosujemy wynik do $-X$,
a dla martyngału średnia jest stała.
\end{proof}

\begin{uwaga}[Dlaczego wymagana jest prawostronna ciągłość filtracji?]
Bez warunku $\calF_{t+}=\calF_t$ granica
$\E[X\mid\calF_r]$ przy $r\downarrow t$ może być
warunkową wartością względem większego $\sigma$-ciała
$\calF_{t+}$, a nie względem $\calF_t$.
Uzyskujemy wtedy proces prawostronnie ciągły, który w chwili
$t$ może wykorzystywać informację niedostępną w $\calF_t$.
To jest powód, dla którego warunki zwykłe są podstawowym
założeniem teorii w czasie ciągłym.
\end{uwaga}


\begin{przyklad}[Dlaczego nie wystarcza lewostronna informacja?]
Niech $\Omega=\{-1,1\}$, obydwa wyniki mają prawdopodobieństwo
$1/2$, a $X(\omega)=\omega$. Przyjmijmy $\calF_0=\{\varnothing,\Omega\}$
oraz $\calF_t=\sigma(X)$ dla $t>0$. Wtedy
$M_t=\E[X\mid\calF_t]$ ma postać $M_0=0$ i $M_t=X$ dla $t>0$.
Jest to martyngał, ale nie posiada modyfikacji prawostronnie ciągłej
w chwili zero. Istotnie, dla wszystkich wymiernych $r>0$
każda modyfikacja musiałaby mieć $\widetilde M_r=X$ poza jednym
wspólnym zbiorem zerowym, a więc prawa granica w zerze wynosiłaby
$X\ne0=\widetilde M_0$ prawie na pewno. Filtracja nie jest
prawostronnie ciągła w zerze, ponieważ
$\calF_{0+}=\bigcap_{t>0}\calF_t=\sigma(X)\supsetneq\calF_0$.
Ten przykład ujawnia, że regularizacja jest twierdzeniem o
\emph{procesie oraz informacji}, a nie wyłącznie o funkcjach czasu.
\end{przyklad}

W dalszej części rozdziału przyjmujemy warunki zwykłe
filtracji, chyba że wyraźnie zaznaczono inaczej.
Adaptowany proces o prawie na pewno c\`adl\`ag
trajektoriach można na mierzalnym zbiorze zerowym
zastąpić zerową trajektorią, zachowując adaptację
i nierozróżnialność. Dlatego w argumentach
trajektoryjnych możemy używać wszędzie c\`adl\`ag
wersji.

\section{Czasy zatrzymania i proces zatrzymany}

\begin{lemat}[Mierzalność wartości procesu zatrzymanego]
Niech $X=(X_t)_{t\ge0}$ będzie rzeczywistym procesem adaptowanym
o wszędzie prawostronnie ciągłych trajektoriach, a $\tau$ czasem zatrzymania.
Dla każdego $t\ge0$ zmienna $X_{t\wedge\tau}$ jest $\calF_t$-mierzalna.
\end{lemat}
\begin{proof}
Ustalmy $t>0$. Przybliżamy $\tau\wedge t$ od góry przez czasy
$\rho_n=2^{-n}\lceil2^n(\tau\wedge t)\rceil\wedge t$, które
przyjmują skończenie wiele wartości z $[0,t]$ i są czasami
zatrzymania. Dla każdego $n$ mamy
$X_{\rho_n}=\sum_{u\in D_n}X_u\chi_{\{\rho_n=u\}}$,
gdzie $D_n$ jest skończonym zbiorem wartości $\rho_n$.
Każdy składnik jest $\calF_t$-mierzalny (ponieważ $u\le t$),
a $\rho_n\downarrow\tau\wedge t$. Prawostronna ciągłość daje
$X_{\rho_n}\to X_{\tau\wedge t}$ punktowo, a więc
ostatnia zmienna jest mierzalna względem $\calF_t$
bez potrzeby uzupełniania filtracji.
Dla $t=0$ jest to adaptowana zmienna $X_0$.
Przy trajektoriach prawostronnie ciągłych jedynie
p.n. ten sam argument działa przy zupełnej filtracji.
\end{proof}


\begin{definicja}[Proces zatrzymany]
Niech $X$ będzie procesem adaptowanym c\`adl\`ag,
a $\tau$ czasem zatrzymania (dopuszczamy wartość $\infty$).
Dla $t\ge0$ definiujemy $X_t^\tau=X_{t\wedge\tau}$.
Gdy $\tau=\infty$, przyjmujemy $t\wedge\tau=t$.
\end{definicja}

\begin{twierdzenie}[Zatrzymanie w ustalonej chwili]
Niech $M$ będzie martyngałem c\`adl\`ag, a $\tau$
czasem zatrzymania. Wtedy $M^\tau$ jest martyngałem.
Jeżeli $X$ jest podmartyngałem (odpowiednio nadmartyngałem)
c\`adl\`ag \emph{klasy $(DL)$}, to $X^\tau$ pozostaje
podmartyngałem (odpowiednio nadmartyngałem).
Tutaj klasa $(DL)$ oznacza, że dla każdego $T<\infty$
rodzina $\{X_\sigma: \sigma\le T,\ \sigma
\text{ jest czasem zatrzymania}\}$ jest jednostajnie
całkowalna.
\end{twierdzenie}
\begin{proof}
Adaptacja wynika z poprzedniego lematu, a trajektorie
zatrzymane są c\`adl\`ag: przed zatrzymaniem
zgadzają się z pierwotną trajektorią, a od chwili
zatrzymania są stałe. Ustalmy $0\le s\le t$ i czas
$\rho=s\vee(t\wedge\tau)$. Ponieważ
$s\le\rho\le t$, twierdzenie o ograniczonym
próbkowaniu opcjonalnym (twierdzenie~\ref{thm:optional-cont})
daje całkowalność $M_\rho,M_{t\wedge\tau},M_{s\wedge\tau}$
oraz $\E[M_\rho\mid\calF_s]=M_s$.
Na zdarzeniu $\{\tau>s\}\in\calF_s$ mamy
$\rho=t\wedge\tau$, a na $\{\tau\le s\}$
proces zatrzymany jest już stały. Dlatego
\[
\begin{aligned}
\E[M_{t\wedge\tau}\mid\calF_s]
&=\chi_{\{\tau\le s\}}M_{s\wedge\tau}
 +\chi_{\{\tau>s\}}\E[M_\rho\mid\calF_s]\\
&=M_{s\wedge\tau}.
\end{aligned}
\]
Dla podmartyngału $X$ klasy $(DL)$ przybliżamy
$\rho$ od góry skończonymi czasami siatkowymi
$\rho_n=2^{-n}\lceil 2^n\rho\rceil\wedge t$,
do siatki włączając $s$ i $t$.
Mamy $s\le\rho_n\le t$. Próbkowanie opcjonalne
w czasie dyskretnym na tej siatce daje
$\E[X_{\rho_n}\mid\calF_s]\ge X_s$.
Po przejściu do granicy otrzymujemy
$\E[X_\rho\mid\calF_s]\ge X_s$.
Stosując ten sam rozkład na zdarzenia
$\{\tau\le s\}$ i $\{\tau>s\}$ co dla martyngału,
dostajemy
$\E[X_{t\wedge\tau}\mid\calF_s]\ge X_{s\wedge\tau}$.
Istotne jest przejście do granicy w $L^1$:
z klasy $(DL)$ otrzymujemy jednostajną całkowalność
wartości w przybliżających czasach zatrzymania,
a prawostronna ciągłość daje ich zbieżność prawie na pewno.
Nierówność przechodzi więc do granicy w warunkowych
wartościach oczekiwanych. Dla nadmartyngału zmieniamy
znak procesu i powtarzamy argument. Samo formalne
podstawienie losowej chwili w definicji podmartyngału
nie jest dowodem.
\end{proof}

\begin{definicja}[Klasy $(D)$ i $(DL)$]
Adaptowany proces c\`adl\`ag jest klasy $(D)$,
jeżeli rodzina wszystkich jego wartości w
ograniczonych czasach zatrzymania jest jednostajnie
całkowalna (bez wspólnego deterministycznego
ograniczenia tych czasów).
Jest klasy $(DL)$, jeżeli dla każdego $T<\infty$
rodzina wartości w czasach zatrzymania $\sigma\le T$
jest jednostajnie całkowalna.
\end{definicja}

\begin{twierdzenie}[Klasy całkowalności martyngału]
\label{thm:mart-cont-class-d}
Każdy martyngał c\`adl\`ag jest klasy $(DL)$.
Dla martyngału c\`adl\`ag równoważne są warunki:
\begin{enumerate}
\item jest klasy $(D)$;
\item rodzina $\{M_t:t\ge0\}$ jest jednostajnie całkowalna;
\item istnieje $X\in L^1$ mierzalna względem
$\calF_\infty=\sigma(\bigcup_{t\ge0}\calF_t)$
taka, że $M_t=\E[X\mid\calF_t]$ p.n. dla każdego $t$.
\end{enumerate}
W ostatnim przypadku dla każdego p.n. skończonego
czasu zatrzymania $\sigma$ zachodzi
$M_\sigma=\E[X\mid\calF_\sigma]$ p.n.
\end{twierdzenie}
\begin{proof}
Dla $\sigma\le T$ ograniczone próbkowanie
opcjonalne (twierdzenie~\ref{thm:optional-cont})
daje $M_\sigma=\E[M_T\mid\calF_\sigma]$.
Lemat~\ref{lem:mart-cont-ui-cond} zapewnia
jednostajną całkowalność na ustalonym horyzoncie.
Z $(D)$ wynika drugi warunek przez wybór czasów
deterministycznych.
Przy drugim warunku dyskretny martyngał $(M_n)$
jest jednostajnie całkowalny, więc zbiega w $L^1$
do $X\in L^1(\calF_\infty)$
(tom I, rozdział~\ref{I-ch:nadmartyngaly}).
Dla każdego ustalonego $t$ i całkowitego $n\ge t$
mamy $M_t=\E[M_n\mid\calF_t]$.
Przejście do granicy w $L^1$ daje trzeci warunek.
Jeśli zachodzi trzeci warunek, dla ograniczonego
$\sigma\le T$ z własności wieży i próbkowania
opcjonalnego
$M_\sigma=\E[M_T\mid\calF_\sigma]
=\E[X\mid\calF_\sigma]$.
Stąd $(D)$ z tego samego lematu.

Dla p.n. skończonego $\sigma$ mamy
$M_{\sigma\wedge n}\to M_\sigma$ p.n.
(ostatecznie $\sigma\wedge n=\sigma$),
a klasa $(D)$ daje również zbieżność w $L^1$.
Dla $A\in\calF_\sigma$ zdarzenie
$A_n=A\cap\{\sigma\le n\}$ należy do
$\calF_{\sigma\wedge n}$, co wynika bezpośrednio
z definicji $\sigma$-ciała w czasie zatrzymania.
Zatem
$\E[M_{\sigma\wedge n}\chi_{A_n}]=\E[X\chi_{A_n}]$.
Przejście do granicy i mierzalność $M_\sigma$
(twierdzenie~\ref{thm:sampled}) identyfikują
żądaną warunkową wartość oczekiwaną.
\end{proof}

\begin{uwaga}[Nie mylić zatrzymania z granicą $M_\tau$]
Powyższe twierdzenie dotyczy procesu $M_{t\wedge\tau}$
przy każdym skończonym $t$. Nie wolno z niego wnioskować
$\E M_\tau=\E M_0$ dla nieograniczonego czasu $\tau$:
przejście $t\to\infty$ wymaga dodatkowego warunku
jednostajnej całkowalności.
\end{uwaga}

\begin{zadanie}[Co odróżnia martyngał od procesu o niemalejącej średniej?]
Niech $W$ będzie standardowym procesem Wienera w jego
filtracji naturalnej. Sprawdzić bez rachunku It\^o, które
z procesów $W_t$, $W_t^2$ i $W_t^2-t$ są martyngałami.

\emph{Rozwiązanie.} Dla $s<t$ zapiszmy
$W_t=W_s+U$, gdzie $U\sim\Normal(0,t-s)$
i $U$ jest niezależne od $\calF_s$.
Stąd $\E[W_t\mid\calF_s]=W_s$.
Po podniesieniu do kwadratu i warunkowaniu otrzymujemy
\[
\E[W_t^2\mid\calF_s]
=W_s^2+2W_s\E U+\E U^2=W_s^2+(t-s).
\]
Proces $W_t^2$ jest zatem podmartyngałem, ale nie
martyngałem. Odejmując od niego deterministyczny
proces $t$, dostajemy
$\E[W_t^2-t\mid\calF_s]=W_s^2-s$.
Wyraz $t$ jest dokładnie kompensacją narastającej
wariancji, którą w kolejnym rozdziale opiszemy za
pomocą nawiasu procesu.
\end{zadanie}

\section{Martyngały lokalne}
\begin{definicja}[Lokalny martyngał i lokalizacja]
Niech $M$ będzie adaptowanym procesem c\`adl\`ag o
wartościach rzeczywistych. Nazywamy go \emph{lokalnym
martyngałem}, jeżeli istnieją czasy zatrzymania
$\tau_k\uparrow\infty$ prawie na pewno, takie że dla
każdego $k$ proces $M^{\tau_k}$ jest martyngałem.
Ciąg $(\tau_k)$ nazywamy \emph{ciągiem lokalizującym}.
W szczególności $M_0\in L^1$, ponieważ
$M^{\tau_k}_0=M_0$. Można zawsze wymagać
$\tau_k\le k$: zastąpienie $\tau_k$ przez
$\tau_k\wedge k$ zachowuje własność martyngałową
na mocy twierdzenia o zatrzymywaniu i nadal daje
ciąg rosnący do nieskończoności p.n.
\end{definicja}

\begin{twierdzenie}[Zatrzymanie ograniczonego procesu lokalnego]
Niech $M$ będzie ciągłym lokalnym martyngałem z $M_0=0$,
a $\tau_K=\inf\{t\ge0:|M_t|\ge K\}$ dla $K>0$.
Wówczas $M^{\tau_K}$ jest martyngałem ograniczonym
przez $K$. Czasy $\tau_K\uparrow\infty$ prawie na pewno,
gdy $K\to\infty$.
\end{twierdzenie}
\begin{proof}
Czasy $\tau_K$ są czasami zatrzymania, ponieważ
dla ciągłej trajektorii i $t\ge0$
$\{\tau_K\le t\}=\{\sup_{0\le u\le t}|M_u|\ge K\}$,
a supremum można obliczyć po chwilach wymiernych
z $[0,t]$ wraz z punktem $t$.
Niech $\sigma_n$ lokalizuje $M$. Proces
$M^{\tau_K\wedge\sigma_n}$ jest martyngałem na mocy
próbkowania opcjonalnego zastosowanego do
$M^{\sigma_n}$: zatrzymujemy go dodatkowo w $\tau_K$.
Ciągłość $M$ daje $|M_{t\wedge\tau_K\wedge\sigma_n}|\le K$.
Dla ustalonego $t$ mamy
$M_{t\wedge\tau_K\wedge\sigma_n}\to M_{t\wedge\tau_K}$
prawie na pewno, a wspólne ograniczenie przez $K$
pozwala przejść do granicy w $L^1$.
Dla $s\le t$ przechodzimy do granicy w warunkowych
równościach martyngałowych, używając kontrakcji warunkowania
w $L^1$. Otrzymujemy martyngał $M^{\tau_K}$.
Na każdym ograniczonym przedziale trajektoria ciągła jest
ograniczona, więc dla dużego $K$ nie przekracza poziomu
$K$ na tym przedziale. To dowodzi $\tau_K\uparrow\infty$.
\end{proof}

\begin{zadanie}[Kompensowany proces Poissona i momenty]
Niech $N$ będzie procesem Poissona o intensywności
$\lambda>0$ oraz $M_t=N_t-\lambda t$. Wykazać,
że $M$ jest martyngałem, i wyznaczyć jego wariancję.

\emph{Rozwiązanie.} Dla $s<t$ przyrost
$N_t-N_s$ ma rozkład Poissona z parametrem
$\lambda(t-s)$ i jest niezależny od $\calF_s$.
Stąd $\E[M_t-M_s\mid\calF_s]=0$. Ponadto
$\Var(M_t)=\Var(N_t)=\lambda t$.
W przeciwieństwie do $W_t$, trajektoria $M$ jest
odcinkami liniowa i ma skoki wielkości $1$.
\end{zadanie}


\begin{twierdzenie}[Nieujemny lokalny martyngał jest nadmartyngałem]
\label{thm:positive-local-supermart}
Niech $M=(M_t)_{t\ge0}$ będzie rzeczywistym lokalnym
martyngałem c\`adl\`ag względem filtracji $(\calF_t)$,
$M_t\ge0$ dla każdego $t$ poza jednym zbiorem zerowym,
oraz $\E M_0<\infty$. Wówczas $M_t\in L^1$ dla wszystkich $t$
i dla $0\le s\le t$ zachodzi
$\E[M_t\mid\calF_s]\le M_s$ prawie na pewno.
W szczególności $\E M_t\le\E M_0$.
\end{twierdzenie}
\begin{proof}
Niech $(\tau_k)$ lokalizuje $M$. Dla każdego $k$ proces
$M^{\tau_k}$ jest martyngałem, toteż
$\E[M_{t\wedge\tau_k}\mid\calF_s]=M_{s\wedge\tau_k}$.
Z warunkowego lematu Fatou dla nieujemnych zmiennych dostajemy
\[
\E[M_t\mid\calF_s]
\le\liminf_{k\to\infty}\E[M_{t\wedge\tau_k}\mid\calF_s]
=\lim_{k\to\infty}M_{s\wedge\tau_k}=M_s.
\]
W tym rachunku lewą stronę początkowo rozumiemy jako rozszerzoną
warunkową wartość oczekiwaną. Niezależnie, z bezwarunkowego
lematu Fatou i $\E M_{t\wedge\tau_k}=\E M_0$ otrzymujemy
$\E M_t\le\E M_0<\infty$ dla każdego $t$.
Wobec tego wszystkie występujące warunkowe wartości
oczekiwane są zwykłymi zmiennymi całkowalnymi, a poprzednia
nierówność dowodzi własności nadmartyngału.
Nierówność może być ostra: lokalizacja zapewnia uczciwość
\emph{przed zatrzymaniem}, ale nie gwarantuje jednostajnej
całkowalności potrzebnej do zachowania średniej po usunięciu
lokalizacji.
\end{proof}

\begin{twierdzenie}[Kryteria dla martyngałów lokalnych]
\label{thm:local-mart-criteria}
Niech $M$ będzie lokalnym martyngałem c\`adl\`ag.
\begin{enumerate}
\item Dla każdego czasu zatrzymania $\rho$
proces $M^\rho$ jest lokalnym martyngałem.
\item Suma i kombinacje liniowe lokalnych
martyngałów są lokalnymi martyngałami.
\item Jeżeli $M$ jest klasy $(DL)$, to jest martyngałem.
Wystarczy również, że dla każdego $T<\infty$
zmienna $\sup_{t\le T}|M_t|$ jest całkowalna.
\item Jeżeli $M\ge0$, to jest martyngałem wtedy
i tylko wtedy, gdy $\E M_t=\E M_0$ dla każdego $t$.
\end{enumerate}
\end{twierdzenie}
\begin{proof}
Niech $(\tau_n)$ lokalizuje $M$.
Zatrzymując martyngał $M^{\tau_n}$ dodatkowo
w $\rho$, dostajemy martyngał
$(M^\rho)^{\tau_n}$, co dowodzi pierwszej tezy.
Dla dwóch procesów z ciągami lokalizującymi
$(\tau_n)$ i $(\sigma_n)$ wspólny ciąg
$(\tau_n\wedge\sigma_n)$ nadal rośnie do
nieskończoności p.n. i lokalizuje ich kombinację
liniową.
Przy klasie $(DL)$ dla ustalonego $t$
rodzina $M_{t\wedge\tau_n}$ jest jednostajnie
całkowalna. Zbieżność p.n. do $M_t$ (ostateczna
równość) daje zbieżność w $L^1$, a stąd
równości martyngałowe przez kontrakcję warunkowania.
Całkowalny supremum na każdym skończonym
horyzoncie daje klasę $(DL)$ przez dominację.
Jeżeli $M\ge0$, to z twierdzenia
\ref{thm:positive-local-supermart}
$\E[M_t\mid\calF_s]\le M_s$.
Przy stałych średnich nieujemna różnica
$M_s-\E[M_t\mid\calF_s]$ ma średnią zero,
więc znika p.n. Odwrotna implikacja jest
bezpośrednia.
\end{proof}

\begin{uwaga}[Różnica między lokalnym martyngałem a martyngałem]
Każdy martyngał c\`adl\`ag jest lokalnym martyngałem
(można wziąć $\tau_k=k$), ale implikacja odwrotna nie jest
ogólnie prawdziwa. Poprzednie twierdzenie pokazuje jedynie,
że nieujemny lokalny martyngał jest nadmartyngałem;
aby stwierdzić równość warunkowych wartości oczekiwanych,
trzeba jeszcze kontrolować brak utraty masy przy lokalizacji.
\end{uwaga}


\begin{twierdzenie}[Odwrotność trójwymiarowego procesu Bessela]
\label{thm:inverse-bessel-strict-local}
Niech $B$ będzie standardowym ruchem Browna
w $\R^3$, $x\ne0$, $r=|x|$ i
$R_t=|x+B_t|$. Proces $R$ nazywamy
trójwymiarowym procesem Bessela rozpoczętym od $r$.
W jego filtracji naturalnej (uzupełnionej)
oraz w filtracji ruchu Browna spełniającej
warunki zwykłe i zachowującej niezależność
przyrostów proces $M_t=1/R_t$ jest dodatnim
ciągłym lokalnym martyngałem, ale nie jest
martyngałem. Dla $t>0$
\[
\E M_t=\frac{2\Phi(r/\sqrt t)-1}{r}
<\frac1r=M_0,
\]
gdzie $\Phi$ jest dystrybuantą standardowego
rozkładu normalnego.
\end{twierdzenie}
\begin{proof}
Z twierdzenia o trafianiu w punkty przez ruch
Browna w rozdziale~\ref{ch:brown-multi}
wynika, że $x+B_t$ prawie na pewno nigdy nie
osiąga zera. Na wspólnym zbiorze zerowym
możemy przyjąć stałą trajektorię $R_t=r$.
Dla $n$ tak dużego, że $1/n<r<n$, połóżmy
\[
\tau_n=\inf\{t\ge0:R_t\notin(1/n,n)\}.
\]
Są to czasy zatrzymania względem obu filtracji,
przez ciągłość trajektorii i twierdzenie
o trafianiu w zbiory domknięte
(twierdzenie~\ref{thm:hitting-closed}).
Na każdym zwartym przedziale $R$ ma dodatnie
minimum i skończone maksimum, więc
$\tau_n\uparrow\infty$ p.n.

Funkcja $h(y)=|y|^{-1}$ jest harmoniczna
na $\R^3\setminus\{0\}$: dla $u=|y|>0$
jej radialny laplasjan wynosi
$h''(u)+2h'(u)/u=2/u^3-2/u^3=0$
(por. twierdzenie~\ref{thm:brown-annulus}).
Wybieramy gładką funkcję o zwartym nośniku
$h_n$ równą $h$ na otoczeniu domkniętego
pierścienia $\{1/n\le |y|\le n\}$,
mnożąc $h$ przez gładką funkcję odcinającą
równą jeden na tym otoczeniu i zero blisko
zera oraz poza dostatecznie dużą kulą.
Lemat~\ref{lem:brown-generator-martingale}
daje ciągły martyngał
\[
h_n(x+B_t)-h_n(x)
-\frac12\int_0^t\Delta h_n(x+B_v)\,\mathrm d\lambda_1(v).
\]
Po zatrzymaniu w $\tau_n$ całka znika,
gdyż $\Delta h_n=0$ na pierścieniu,
a pozostała różnica wynosi
$M_{t\wedge\tau_n}-1/r$.
Zatem $M^{\tau_n}$ jest martyngałem
w filtracji ruchu Browna. W mniejszej filtracji
$R$ pozostaje nim przez własność wieży
i adaptację. Dowód lokalności nie wymaga
rachunku It\^o.

Obliczymy teraz średnią bez lokalizacji.
Dzięki symetrii obrotowej możemy przyjąć
$x=(0,0,r)$. Dla $t>0$ gęstość $x+B_t$
w punkcie $y$ wynosi
$(2\pi t)^{-3/2}\exp(-|y-x|^2/(2t))$.
We współrzędnych sferycznych
$\rho=|y|>0$, $\alpha$ jest kątem względem osi
$x$, a element objętości ma postać
$\rho^2\sin\alpha\,\mathrm d\lambda_3(\rho,\alpha,\beta)$.
Po całkowaniu po kącie azymutalnym $\beta$
oraz podstawieniu $v=\cos\alpha$ otrzymujemy
\[
\begin{aligned}
\E M_t
&=\frac{2\pi}{(2\pi t)^{3/2}}\int_0^\infty
\rho e^{-(\rho^2+r^2)/(2t)}
\left(\int_{-1}^1 e^{\rho rv/t}\,\mathrm d\lambda_1(v)\right)
\,\mathrm d\lambda_1(\rho)\\
&=\frac1{r\sqrt{2\pi t}}\int_0^\infty
\left(e^{-(\rho-r)^2/(2t)}-
e^{-(\rho+r)^2/(2t)}\right)
\,\mathrm d\lambda_1(\rho)\\
&=\frac1r\left(\Phi(r/\sqrt t)-\Phi(-r/\sqrt t)\right).
\end{aligned}
\]
W drugiej równości całkujemy funkcję wykładniczą
po $v$, w ostatniej podstawiamy
$(\rho-r)/\sqrt t$ oraz $(\rho+r)/\sqrt t$.
Symetria rozkładu normalnego daje żądany wzór.
Dla każdego $t>0$ średnia jest ściśle mniejsza
niż $1/r$, więc proces nie jest martyngałem.
Jest to \emph{ścisły lokalny martyngał},
w odróżnieniu od prawdziwego martyngału,
który traci średnią dopiero w granicy
$t\to\infty$.
\end{proof}

\section{Wykładnik gaussowski: martyngał, który traci średnią w granicy}
\begin{twierdzenie}[Wykładnik procesu Wienera]
Niech $W$ będzie standardowym procesem Wienera, a
$\theta\in\R$. Wtedy
$Z_t=\exp(\theta W_t-\tfrac12\theta^2t)$
jest dodatnim, ciągłym martyngałem o $\E Z_t=1$
dla każdego $t\ge0$. Gdy $\theta\ne0$,
$Z_t\to0$ prawie na pewno przy $t\to\infty$.
W szczególności rodzina $\{Z_t:t\ge0\}$ nie jest
jednostajnie całkowalna. Jest to konkretny przykład
martyngału ograniczonego w $L^1$, który nie zbiega w $L^1$.
\end{twierdzenie}
\begin{proof}
Niech $s<t$. Z niezależności przyrostów
$W_t-W_s\sim\Normal(0,t-s)$ i z gaussowskiej
funkcji tworzącej momenty mamy
\[
\E[Z_t\mid\calF_s]
=Z_s\E\exp\bigl(\theta(W_t-W_s)-\tfrac12\theta^2(t-s)\bigr)
=Z_s.
\]
Dla liczb całkowitych $n$, przyrosty
$W_n-W_{n-1}$ są niezależne i mają ten sam rozkład
$\Normal(0,1)$; mocne prawo wielkich liczb daje
$W_n/n\to0$ prawie na pewno.
Dla $u>0$ zasada odbicia (twierdzenie~\ref{thm:wiener-reflection})
i symetria dają
$\Prob(\sup_{0\le s\le1}|W_s|>u)\le
4(1-\Phi(u))\le C e^{-u^2/2}$ dla $u\ge1$.
Dlatego dla każdego $\varepsilon>0$ suma
$\sum_n\Prob(\sup_{n\le t\le n+1}|W_t-W_n|>\varepsilon n)$
jest skończona. Lemat Borela--Cantellego pokazuje, że
$\sup_{n\le t\le n+1}|W_t-W_n|/n\to0$ prawie na pewno.
Stąd $W_t/t\to0$ dla $t\to\infty$, a więc
$t^{-1}\log Z_t=\theta W_t/t-\theta^2/2\to-\theta^2/2<0$.
To wymusza $Z_t\to0$ prawie na pewno.
Gdyby rodzina $(Z_t)$ była jednostajnie całkowalna,
zbieżność prawie na pewno pociągałaby zbieżność w $L^1$,
a zatem $\E Z_t\to0$. Przeczy to $\E Z_t=1$.
\end{proof}

\section{Przypadek zatrzymania na pierwszym trafieniu}
\begin{zadanie}[Dlaczego $\E M_\tau$ może różnić się od $\E M_0$?]
Niech $W$ będzie procesem Wienera, a
$\tau=\inf\{t\ge0:W_t=1\}$. Wykazać, że
$\tau<\infty$ prawie na pewno, lecz
$\E W_\tau=1\ne0=\E W_0$.

\emph{Rozwiązanie.} Z zasady odbicia (twierdzenie~\ref{thm:wiener-reflection})
$\Prob(\tau\le t)=2(1-\Phi(1/\sqrt t))\to1$.
Z ciągłości trajektorii $W_\tau=1$ na $\{\tau<\infty\}$.
Dla każdego skończonego $T$ proces $W_{t\wedge\tau}$
jest martyngałem i $\E W_{T\wedge\tau}=0$.
Nie można jednak przejść do granicy pod wartością oczekiwaną:
rodzina $\{W_{T\wedge\tau}:T\ge0\}$ nie jest
jednostajnie całkowalna, skoro granica prawie na pewno
ma wartość oczekiwaną $1$, a wszystkie zmienne mają
średnią $0$. To konkretny przykład nieuprawnionego
stosowania próbkowania opcjonalnego.
\end{zadanie}


% ================================================================
% ROZDZIAŁ 12. PRÓBKOWANIE, NIERÓWNOŚCI, ZBIEŻNOŚĆ
% ================================================================
